\chapter{Implementation: Multipath Real-Time Video Testbed}
\label{ch:p2impl}

Realizing the Part~II design requires a testbed that is faithful enough to expose real multipath and
congestion behavior yet fast enough for reinforcement-learning rollouts. We resolve this tension with a
single data-plane abstraction served by two interchangeable backends (\cref{sec:p2impl:dataplane}): a
packet-level NS-3 simulator and a fast analytical mock. \Cref{sec:p2impl:rt} describes the real-time
video and multipath model, \cref{sec:p2impl:agents} the agent implementation and the VMAF surrogate, and
\cref{sec:p2impl:gaps} the known gaps that scope the current results.

\section{Data-Plane Abstraction and Two Backends}
\label{sec:p2impl:dataplane}

Both backends implement one abstract \texttt{DataPlane} interface, so the RL code, observations, and
rewards are written once and run unchanged against either. This lets us iterate quickly on the mock and
validate on the higher-fidelity simulator without divergence.

\begin{description}
  \item[NS-3 backend.] A packet-level simulation in which the connection's several paths are modeled as
  independent subflows over a topology of links, each with its own capacity, propagation delay, and
  active-queue management (FqCoDel). Realistic contention comes from OnOff UDP cross-traffic with
  exponentially distributed on/off periods, phase-shifted per path so the paths do not degrade in
  lockstep. Each video frame's bytes are striped across the subflows according to the transport agent's
  split, and delivery is tracked by a byte-watermark scheme -- avoiding per-packet tagging -- from which
  per-frame latency, jitter, and deadline-miss loss are measured. Python and NS-3 communicate through an
  ns3-ai shared-memory bridge.
  \item[Mock backend.] A pure-Python, per-seed-deterministic model in which each path has a capacity
  envelope (a smooth sinusoidal component plus bursty cross-traffic and noise) and a standing-queue model
  that produces bufferbloat when a path is overdriven. It runs orders of magnitude faster than NS-3 and is
  used for rapid training and for scenarios the NS-3 body does not yet emit.
\end{description}

\paragraph{Non-stationary dynamics.} The mock backend can optionally inject the volatility that makes
multipath scheduling matter: \emph{path churn} (paths appear and disappear, and bytes sent to a dead path
are lost), \emph{regime shifts} (abrupt changes in the best path's capacity), \emph{congestion bursts}
(transient per-path capacity collapses), and \emph{correlated failures} (shared-bottleneck groups that
degrade together). These are off by default and are the conditions under which \cref{ch:p2eval} finds the
learned split decisive.

\section{Real-Time Video and Multipath Model}
\label{sec:p2impl:rt}

The video source mirrors a conferencing encoder: a fixed frame rate (30\,fps by default), variable
frame sizes (with periodic larger I-frames and per-frame jitter), no playback buffer, and a deadline
after which a late frame is discarded as lost. The aggregate target bitrate ranges over
\SIrange{300}{6000}{\kilo\bit\per\second}. The canonical topology is deliberately \emph{non-dominant}: no
single path can carry full-quality video, and the aggregate usable capacity sits below the quality knee,
so the system must genuinely combine paths and trade rate against delivery. For example, a four-path
configuration mixes paths of differing capacity and RTT -- including a higher-capacity but high-latency
``trap'' path -- under heavy, path-specific cross-traffic, so that neither the widest nor the
lowest-latency path is right on its own.

\section{Agents and QoE Surrogate}
\label{sec:p2impl:agents}

\paragraph{Agents.} Both agents are soft actor-critic learners with twin critics and automatic
temperature tuning. The transport agent has two variants: a \emph{flat} version whose actor and critic
consume a fixed-length concatenation of the global and per-path features (a fixed path count), and a
\emph{scoring} version implementing the permutation-equivariant, mask-aware design of
\cref{sec:p2design:scoring} with a shared per-path actor and a Deep Sets masked-mean-pool critic. The
hierarchical training loop runs the two agents at their respective cadences, storing each agent's
transitions in its own replay buffer and crediting the application agent over the window of frames its
bitrate governed.

\paragraph{VMAF surrogate.} The quality term of \cref{eq:p2qoe} needs a mapping from an encoding (and
network conditions) to perceptual quality. Two are supported: a default bitrate-only logarithmic curve fit
through published quality anchors, and a learned surrogate that interpolates a
quality-of-service$\rightarrow$VMAF table over bitrate, loss, delay, and jitter (produced by a sibling
data-generation project). When the learned surrogate is active, the explicit loss penalty in
\cref{eq:p2qoe} is dropped to avoid double-counting the loss it already reflects.

\section{Known Gaps}
\label{sec:p2impl:gaps}

The current testbed has three limitations that scope the results of \cref{ch:p2eval} and define concrete
follow-up work.

\begin{itemize}
  \item \textbf{Dynamics parity.} The non-stationary dynamics (churn, regime shifts, bursts) are
  currently emitted only by the mock backend; porting them into the NS-3 body -- adding per-path liveness
  to the simulation -- is required before the packet-level and analytical results can be compared on equal
  footing. \gap{report NS-3 vs.\ mock parity once the C++ dynamics land.}
  \item \textbf{Transport realism.} Subflows are modeled with TCP-like congestion control rather than true
  Multipath QUIC; this preserves congestion, queueing, and loss dynamics but masks some network loss as
  latency. Swapping in a real QUIC module is future work.
  \item \textbf{Surrogate coverage.} The learned VMAF surrogate's loss, delay, and jitter axes are
  currently under-sampled, so the quality term is effectively bitrate-dominated; a richer surrogate would
  activate those axes and let quality respond to delivery conditions directly.
\end{itemize}
